\section{Plan de tests}
\label{ch:tests}

La stratégie de tests couvre quatre niveaux : tests \textbf{unitaires} (composants isolés),
tests d'\textbf{intégration} (parcours traversant plusieurs couches), tests de
\textbf{sécurité} (rejouant les failles corrigées) et tests de \textbf{non-régression}. Le
tableau~\ref{tab:plantests} présente le plan, dérivé des fonctionnalités et des \emph{user
stories} de la section~\ref{sec:acteurs}.

\begin{table}[h]
\centering
\renewcommand{\arraystretch}{1.3}
\small
\begin{tabularx}{\textwidth}{@{}L{3.6cm} L{2.6cm} X@{}}
\toprule
\textbf{Fonctionnalité} & \textbf{Type de test} & \textbf{Objet du test} \\
\midrule
Connexion & Unitaire & \texttt{UserValidator::login} (champs vides, email inconnu, mauvais mot de passe) \\
Échappement & Unitaire & \texttt{Utils::e()} neutralise \texttt{<script>} \\
Accès données & Unitaire & \texttt{PostModel} : création, lecture, suppression \\
Routage & Unitaire & \texttt{Router} : correspondance d'URL et paramètres \\
Publier un message & Intégration & Parcours complet : \texttt{fetch} $\to$ contrôleur $\to$ modèle $\to$ base \\
Injection SQL & Sécurité & Recherche avec charge utile : la requête préparée neutralise \\
XSS & Sécurité & Message \texttt{<script>} : rendu échappé, non exécuté \\
CSRF & Sécurité & POST sans jeton $\to$ \texttt{403} \\
IDOR & Sécurité & Suppression d'un message d'autrui $\to$ \texttt{403} \\
\bottomrule
\end{tabularx}
\caption{Plan de tests de \projet{}.}
\label{tab:plantests}
\end{table}

\begin{todo}{Implémentation et exécution des tests (PHPUnit)}
Cette section doit être complétée par l'implémentation réelle — c'est la priorité du projet, la
compétence C9 étant \textbf{obligatoire}. À produire :
\begin{enumerate}
  \item installer \textbf{PHPUnit} (dépendance de développement Composer) ;
  \item créer un \textbf{environnement de test} : conteneur PostgreSQL dédié (ou SQLite en
        mémoire via un adaptateur) ;
  \item écrire les \textbf{tests unitaires} listés ci-dessus et capturer la sortie verte ;
  \item écrire les \textbf{quatre tests de sécurité} (injection SQL, XSS, CSRF, IDOR) ;
  \item insérer une \textbf{capture du rapport d'exécution} (nombre de tests, assertions,
        verts / rouges) ;
  \item renseigner la colonne « résultat obtenu » du jeu d'essai (section suivante).
\end{enumerate}
Code de test, sorties complètes et jeux de données iront en annexe~\ref{ann:tests}.
\end{todo}

\section{Jeu d'essai de la fonctionnalité la plus représentative}
\label{sec:jeuessai}

La fonctionnalité retenue est « \textbf{Publier un message} » : elle traverse toutes les couches
(interface, contrôleur, validation, téléversement, modèle, SQL, réponse, rendu) et concentre les
protections de sécurité. Le tableau~\ref{tab:jeuessai} en présente le jeu d'essai.

\begin{table}[h]
\centering
\renewcommand{\arraystretch}{1.3}
\small
\begin{tabularx}{\textwidth}{@{}L{2.6cm} L{3.2cm} X L{2.2cm}@{}}
\toprule
\textbf{Cas} & \textbf{Donnée d'entrée} & \textbf{Résultat attendu} & \textbf{Obtenu} \\
\midrule
Nominal & Contenu « Bonjour » & Message créé, \texttt{id} retourné, HTTP 200 & \textit{(à exécuter)} \\
Contenu vide & \texttt{""} & Refus « contenu vide », pas d'insertion & \textit{(à exécuter)} \\
Image trop lourde & Fichier > limite & Refus « image trop lourde » & \textit{(à exécuter)} \\
Extension interdite & Fichier \texttt{.php} & Refus « type de fichier non autorisé » & \textit{(à exécuter)} \\
Charge XSS & \texttt{<script>\dots} & Stocké brut, \textbf{rendu échappé} & \textit{(à exécuter)} \\
Non authentifié & Aucune session & Redirection / \texttt{403} (AuthMiddleware) & \textit{(à exécuter)} \\
Sans jeton CSRF & POST sans jeton & \texttt{403} (CsrfMiddleware) & \textit{(à exécuter)} \\
\bottomrule
\end{tabularx}
\caption{Jeu d'essai de la fonctionnalité « Publier un message ».}
\label{tab:jeuessai}
\end{table}

\begin{todo}{Colonne « résultat obtenu » et analyse des écarts}
Après exécution, renseigner la colonne « Obtenu » pour chaque cas et ajouter, sous le tableau,
l'\textbf{analyse des écarts} éventuels (cas où l'obtenu diffère de l'attendu, cause, correctif).
Le référentiel exige cette analyse, même en l'absence d'écart (« analyse des écarts éventuels »).
\end{todo}

\section{Préparation et documentation du déploiement}
\label{sec:deploiement}

La mise en production s'appuie sur la pile \file{docker-compose.prod.yml} : Traefik termine le
HTTPS (Let's Encrypt), redirige HTTP vers HTTPS et pose l'en-tête HSTS ; le service
\texttt{migrate} applique les migrations \emph{avant} le démarrage de l'application ; PostgreSQL
n'est pas exposé hors du réseau interne.

\begin{lstlisting}[style=yaml, caption={Routage HTTPS et HSTS en production — \file{docker-compose.prod.yml}}]
labels:
  - traefik.http.routers.solage.entrypoints=websecure
  - traefik.http.routers.solage.tls.certresolver=le
  - traefik.http.middlewares.solage-hsts.headers.stsSeconds=31536000
  - traefik.http.middlewares.solage-hsts.headers.stsIncludeSubdomains=true
  - traefik.http.routers.solage.middlewares=solage-hsts
\end{lstlisting}

\begin{todo}{Procédure de déploiement (\file{DEPLOYMENT.md})}
Rédiger une procédure de déploiement documentée (critère C10) :
\begin{itemize}
  \item \textbf{pré-requis} (serveur, Docker, variables d'environnement, nom de domaine, DNS) ;
  \item \textbf{étapes} de mise en production (build de l'image, \texttt{docker compose up},
        migrations, vérification) ;
  \item procédure de \textbf{retour arrière} (\emph{rollback}) en cas d'incident ;
  \item description des \textbf{environnements} (développement local, pré-production, production)
        et de leurs différences.
\end{itemize}
\end{todo}

\section{Contribution à la mise en production (DevOps)}
\label{sec:devops}

Plusieurs briques DevOps sont déjà en place : conteneurisation complète (\texttt{docker
compose}), images reproductibles, migrations \textbf{idempotentes} rejouables sans risque, et
journalisation structurée adaptée à une collecte par l'orchestrateur. Il manque l'automatisation
par intégration continue.

\begin{todo}{Pipeline d'intégration continue (CI/CD)}
Mettre en place un pipeline \textbf{GitHub Actions} déclenché à chaque \emph{push}, et y joindre
une capture d'un rapport d'exécution. Étapes cibles : analyse statique (PHPStan), tests
(PHPUnit), \emph{lint} JavaScript (ESLint), construction de l'image Docker. Squelette de départ à
adapter :
\end{todo}

\begin{lstlisting}[style=yaml, caption={Squelette de pipeline CI proposé — \file{.github/workflows/ci.yml} (à implémenter)}]
name: CI
on: [push]
jobs:
  build-and-test:
    runs-on: ubuntu-latest
    steps:
      - uses: actions/checkout@v4
      - uses: shivammathur/setup-php@v2
        with: { php-version: '8.3' }
      - run: composer install --no-interaction
      - run: vendor/bin/phpstan analyse src modules   # analyse statique
      - run: vendor/bin/phpunit                        # tests unitaires + securite
      - run: docker build -t solage-app .              # construction du livrable
\end{lstlisting}

\section{Veille de sécurité}
\label{sec:veille}

\paragraph{Sources suivies.} La veille s'appuie sur l'OWASP (Top~10 et guides de test), les
publications de l'ANSSI, les bulletins de sécurité de PHP et les bases de vulnérabilités (CVE).

\paragraph{Le déclic IDOR.} Un article sur la fuite de données ayant touché l'URSSAF (accès aux
données d'autres assurés en modifiant un identifiant dans l'URL — une faille \textbf{IDOR}) a
déclenché un audit de \projet{}. Il a révélé que les routes \texttt{/edituser/\{id\}},
\texttt{/api/users/delete} et \texttt{/api/posts/delete} vérifiaient l'\emph{authentification}
(\texttt{AuthMiddleware}) mais pas la \emph{propriété} de la ressource : un utilisateur
authentifié pouvait modifier ou supprimer le contenu d'autrui en changeant l'identifiant. Le
correctif (contrôle « propriétaire ou administrateur » + \texttt{403} + journalisation) est
décrit en section~\ref{sec:metier}.

\begin{todo}{Journal de veille}
Compléter par un \textbf{journal de veille} daté : pour chaque entrée, la source, la date, ce qui
a été appris, la vulnérabilité éventuellement détectée dans \projet{} et le correctif apporté.
Viser au moins deux à trois entrées en plus de l'IDOR.
\end{todo}
